\section{Introduction}
Diffusion probabilistic models (DPMs)~\citep{sohl2015deep,ho2020denoising,song2020score} are emerging powerful generative models with promising performance on many tasks, such as image generation~\citep{dhariwal2021diffusion,meng2021sdedit}, video generation~\citep{ho2022video}, text-to-image generation~\citep{ramesh2022hierarchical}, speech synthesis~\citep{chen2020wavegrad,chen2021wavegrad} and lossless compression~\citep{kingma2021variational}.
DPMs are defined by discrete-time random processes~\citep{sohl2015deep,ho2020denoising} or continuous-time stochastic differential equations (SDEs)~\citep{song2020score}, which learn to gradually remove the noise added to the data points.
Compared with the widely-used generative adversarial networks (GANs)~\citep{goodfellow2014generative} and variational auto-encoders (VAEs)~\citep{kingma2013auto}, DPMs can not only compute exact likelihood~\citep{song2020score}, but also achieve even better sample quality for image generation~\citep{dhariwal2021diffusion}. 
However, to obtain high-quality samples, DPMs usually need hundreds or thousands of sequential steps of large neural network evaluations, thereby resulting in a much slower sampling speed than the single-step GANs or VAEs. Such inefficiency is becoming a critical bottleneck for the adoption of DPMs in downstream tasks, leading to an urgent request to design fast samplers for DPMs.




Existing fast samplers for DPMs can be divided into two categories.
The first category includes knowledge distillation~\citep{salimans2022progressive,luhman2021knowledge} and noise level or sample trajectory learning~\citep{san2021noise,nichol2021improved,lam2021bilateral,watson2021learning}. Such methods require a possibly expensive training stage before they can be used for efficient sampling. Furthermore, their applicability and flexibility might be limited. It might require nontrivial effort to adapt the method to different models, datasets, and number of sampling steps. 
The second category consists of training-free~\citep{song2020denoising,jolicoeur2021gotta,bao2022analytic} samplers, which are suitable for all pre-trained DPMs in a simple plug-and-play manner. 
Training-free samplers include adopting implicit~\citep{song2020denoising} or analytical~\citep{bao2022analytic} generation process, advanced differential equation (DE) solvers~\citep{song2020score,jolicoeur2021gotta,liu2022pseudo,popov2021diffusion,tachibana2021taylor} and dynamic programming~\citep{watson2021learning}.
However, these methods still require $\sim$ 50 function evaluations~\citep{bao2022analytic} to generate high-quality samples (comparable to those generated by plain samplers in about 1000 function evaluations), thereby are still time-consuming. 












In this work, we bring the efficiency of training-free samplers to a new level to produce high-quality samples in the ``\textit{few-step sampling}'' regime, where the sampling can be done within around 10 steps of sequential function evaluations. We tackle the alternative problem of sampling from DPMs as solving the corresponding diffusion ordinary differential equations (ODEs) of DPMs, and carefully examine the structure of diffusion ODEs. 
Diffusion ODEs have a semi-linear structure --- they consist of a linear function of the data variable and a  nonlinear function  parameterized by neural networks. 
Such structure is omitted in previous training-free samplers~\cite{song2020score,jolicoeur2021gotta}, which directly use black-box DE solvers. 
To utilize the semi-linear structure, we derive an exact formulation of the solutions of diffusion ODEs by analytically computing the linear part of the solutions, avoiding the corresponding discretization error. 
Furthermore, by applying change-of-variable, the solutions can be equivalently simplified to an exponentially weighted integral of the neural network. Such integral is very special and can be efficiently approximated by the numerical methods for exponential integrators~\citep{hochbruck2010exponential}. 








Based on our formulation of solutions, we propose \textit{DPM-Solver}, a fast dedicated solver for diffusion ODEs by approximating the above integral. Specifically, we propose first-order, second-order and third-order versions of DPM-Solver with convergence order guarantees.
We further propose an adaptive step size schedule for DPM-Solver.
In general, DPM-Solver is applicable to both continuous-time and discrete-time DPMs, and also conditional sampling with classifier guidance~\citep{dhariwal2021diffusion}.
Fig.~\ref{fig:intro} demonstrates the speedup performance of a Denoising Diffusion Implicit Models (DDIM)~\citep{song2020denoising} baseline and DPM-Solver, which shows that DPM-Solver can generate high-quality samples with as few as 10 function evaluations and is much faster than DDIM on the ImageNet 256x256 dataset~\citep{deng2009imagenet}. Our additional experimental results show that DPM-Solver can greatly improve the sampling speed of both discrete-time and continuous-time DPMs, and it can achieve excellent sample quality in around 10 function evaluations, which is much faster than all previous training-free samplers of DPMs.


\begin{figure}[t]
	\begin{minipage}{0.18\linewidth}
		\centering
		  %
			\includegraphics[width=\linewidth]{figures/ddim_intro_1_10.png}\\
			\includegraphics[width=\linewidth]{figures/ddim_intro_2_10.png}\\
			\small{NFE = 10} \\
	\end{minipage}
	\hspace{-0.15cm}
	\begin{minipage}{0.18\linewidth}
		\centering
		  %
			\includegraphics[width=\linewidth]{figures/ddim_intro_1_15.png}\\
			\includegraphics[width=\linewidth]{figures/ddim_intro_2_15.png}\\
		    \small{NFE = 15} \\
	\end{minipage}
	\hspace{-0.15cm}
	\begin{minipage}{0.18\linewidth}
		\centering
		  %
			\includegraphics[width=\linewidth]{figures/ddim_intro_1_20.png}\\
			\includegraphics[width=\linewidth]{figures/ddim_intro_2_20.png}\\
		    \small{NFE = 20} \\
	\end{minipage}
	\hspace{-0.15cm}
	\begin{minipage}{0.18\linewidth}
		\centering
		  %
			\includegraphics[width=\linewidth]{figures/ddim_intro_1_100.png}\\
			\includegraphics[width=\linewidth]{figures/ddim_intro_2_100.png}\\
		    \small{NFE = 100} \\
	\end{minipage}
	\hspace{-0.15cm}
	\hspace{0.05\linewidth}
	\begin{minipage}{0.18\linewidth}
		\centering
		  %
			\includegraphics[width=\linewidth]{figures/dpm_solver_intro_1_10.png}\\
			\includegraphics[width=\linewidth]{figures/dpm_solver_intro_2_10.png}\\
		    \small{NFE = 10} \\
	\end{minipage}
	
	\vspace{0.3cm}
	\begin{minipage}{\linewidth}
    \hspace{0.28\linewidth}
    \small{(a) DDIM~\citep{song2020denoising}}
    \hspace{0.34\linewidth}
    \small{(b) DPM-Solver (ours)}
	\end{minipage}
	
	\caption{\small{Samples by DDIM~\citep{song2020denoising} with 10, 15, 20, 100 number of function evaluations (NFE), and DPM-Solver (ours) with only 10 NFE, using the pre-trained DPMs on ImageNet 256$\times$256 with classifier guidance~\citep{dhariwal2021diffusion}.}
	\label{fig:intro}
}
\vspace{-.2in}
\end{figure}




